\documentclass[conference]{IEEEtran}
\usepackage[utf8]{inputenc}
\usepackage{amsmath}
\usepackage{amssymb}
\usepackage{graphicx}
\usepackage{cite}
\usepackage{algorithm}
\usepackage{algpseudocode}
\usepackage{tikz}
\usetikzlibrary{arrows.meta,positioning,shapes.geometric,shapes.misc,fit,backgrounds}
\usepackage{float}

\tikzset{
  block/.style={rectangle, draw, rounded corners, align=center, minimum height=1.2cm, minimum width=3.2cm},
  process/.style={rectangle, draw, align=center, minimum height=1.0cm, minimum width=3.3cm},
  io/.style={trapezium, trapezium left angle=75, trapezium right angle=105, draw, align=center, minimum height=1.0cm, minimum width=3.3cm},
  decision/.style={diamond, draw, aspect=2.0, align=center, inner sep=1pt},
  data/.style={cylinder, draw, shape border rotate=90, aspect=0.25, align=center, minimum height=1.1cm, minimum width=2.0cm},
  line/.style={-{Latex[length=2.2mm,width=2.2mm]}, thick},
  dashedline/.style={-{Latex[length=2mm,width=2mm]}, dashed}
}

\title{Comparative Analysis of RSSI and CSI-Based Localization for Hidden IoT Devices in Cluttered Indoor Environments Using Hybrid Filtering Algorithms}

\author{
    \IEEEauthorblockN{Mac Dylan Philippe M. Balagtas$^{[1]}$, Niño Angelo U. Balagtas$^{[1]}$, Alec Xavier B. Que Esteves$^{[1]}$, \\ Dr. Luis Gerardo Jr. S. Cañete$^{[2]}$, Dr. Philip Virgil B. Astillo$^{[2]}$}
    \IEEEauthorblockA{\textit{$^{[1]}$Undergraduate Student, Department of Computer Engineering, University of San Carlos, Talamban, Cebu City, Philippines}}
    \IEEEauthorblockA{\textit{$^{[2]}$Faculty Member, Department of Computer Engineering, University of San Carlos, Talamban, Cebu City, Philippines}}
    \IEEEauthorblockA{ 07501508@usc.edu.ph, 24105130@usc.edu.ph, 22102490@usc.edu.ph,}
    \IEEEauthorblockA{ lscanete@usc.edu.ph, pvbastillo@usc.edu.ph}
}

\begin{document}

\maketitle

\begin{abstract}
The proliferation of unauthorized or hidden IoT devices in sensitive indoor environments presents a growing challenge for privacy and security. While Wi-Fi-based Indoor Positioning Systems (IPS) have matured, localizing ``hidden'' stationary nodes in cluttered spaces remains difficult due to severe multipath propagation and signal shadowing. Traditional localization methods rely on Received Signal Strength Indicator (RSSI), which exhibits high variance in indoor settings. This research proposes a robust Real-Time Location System (RTLS) that fuses RSSI with fine-grained Channel State Information (CSI) to enhance localization accuracy for non-cooperative IoT devices. Utilizing a distributed architecture of ESP32 edge nodes and a laptop gateway, the study introduces an adaptive hybrid filter to mitigate indoor RF noise. Experimental results, conducted in a semi-structured classroom environment at the University of San Carlos, demonstrate that the hybrid approach significantly reduces Mean Absolute Error (MAE) compared with baseline RSSI methods, providing a low-cost framework for indoor security and asset tracking.\smallskip
\end{abstract}

\begin{IEEEkeywords}
Indoor Localization, CSI, RSSI, ESP32, Multipath Propagation, Hidden IoT Detection, Physical Layer Security.
\end{IEEEkeywords}

\section{Introduction}
The rapid proliferation of the Internet of Things (IoT) has led to an increase in the deployment of small, Wi-Fi-enabled devices for various applications, ranging from smart home automation to industrial monitoring. However, this ubiquity introduces significant physical-layer security risks, such as the placement of unauthorized ``hidden'' devices—including rogue access points or illicit surveillance nodes—within private or educational spaces. Detecting and localizing these non-cooperative devices is critical for maintaining privacy and preventing data breaches \cite{ren2022lumos,ju2025hidden}.\smallskip

Current localization systems primarily rely on Received Signal Strength Indicator (RSSI) due to its universal support and low computational cost \cite{zaruba2007singleap}. However, in cluttered indoor settings, RSSI is highly susceptible to extreme multipath propagation, leading to localization errors that can exceed several meters \cite{yang2013rssi2csi,wu2012csi}. Comparative outdoor LoRa IoT measurements likewise report stronger robustness for CSI-derived features than for RSSI under spatial and temporal variability \cite{ghany2020robustness}, supporting CSI--RSSI tradeoff analysis even though that campaign differs from indoor Wi-Fi clutter. To address these limitations, this study explores the fusion of RSSI and Channel State Information (CSI) to create a robust Real-Time Location System (RTLS), following established CSI-based indoor localization literature \cite{yang2013rssi2csi,wu2012csi,li2024efficient}. By utilizing low-cost ESP32 microcontrollers as passive sniffers, this research seeks to bridge the gap between high-end academic localization theory and accessible, real-world security hardware.\smallskip

\subsection{Problem Background}
On a global scale, the explosion of IoT connectivity has outpaced the development of localized security frameworks. Internationally, the difficulty of precisely localizing devices that do not ``cooperate'' with network infrastructure remains a significant barrier to digital privacy. While high-accuracy localization has been achieved in laboratory settings, most current solutions are either too expensive for mass deployment or rely on active network handshakes that hidden devices are designed to avoid.\smallskip

In the urban context of Cebu City, the challenge becomes more acute. In dense, semi-structured environments such as the University of San Carlos--Talamban Campus, traditional long-range localization methods such as GPS are rendered ineffective by indoor signal attenuation and urban canyon effects \cite{sensors2022indoor3d,xie2015urbancanyon}. This forces a reliance on Wi-Fi-based indoor positioning. However, past research in this area typically requires a user to walk around the environment to ``map'' signal fingerprints---a requirement that is impractical for the stealthy detection of stationary hidden devices \cite{wu2012csi,yang2013rssi2csi}.\smallskip

At the most granular level—within the individual classrooms of the USC-TC Bunzel Building—the problem reaches its technical peak. These spaces act as an ``RF Hall of Mirrors,'' characterized by extreme multipath propagation where signals reflect off concrete walls and metal furniture. In these specific micro-environments, a hidden device can easily ``mask'' its location within signal nulls. Most existing academic frameworks utilize specialized, expensive network cards (like the Intel 5300) to solve this, leaving a gap for research that utilizes ubiquitous, low-cost hardware in a static, mesh-based architecture \cite{yang2013rssi2csi,wu2012csi}. This research addresses this localization gap by leveraging the fine-grained physical layer data of Wi-Fi signals to overcome the limitations of both RSSI-only and movement-dependent systems.\smallskip

\subsection{Statement of the Problem}
Wi-Fi-based indoor localization for stationary, hidden IoT devices faces a technical trilemma of precision, deployment autonomy, and computational cost. Specifically, this research addresses the following identified problems:

\begin{enumerate}
    \item \textbf{The Reliability Crisis of RSSI in Extreme Multipath:} In cluttered indoor environments like the USC-TC Bunzel Building, signals do not follow a predictable path. Concrete walls and metal furniture create an ``RF Hall of Mirrors'' where constructive and destructive interference causes Received Signal Strength Indicator (RSSI) values to fluctuate wildly. This leads to localization errors of several meters, rendering standard trilateration ineffective for finding small, hidden devices.
    
    \item \textbf{The Mobility Dependency of Existing Frameworks:} A significant portion of high-accuracy indoor positioning research relies on dynamic fingerprinting or user-led site surveys. Such methods require a human operator to physically walk around a room to map signal variations. This is impractical for security applications where stealth and speed are required to localize non-cooperative, stationary hidden nodes without prior environmental mapping.
    
    \item \textbf{Computational Bottlenecks at the Resource-Constrained Edge:} While Channel State Information (CSI) provides the granularity needed to overcome multipath noise, its raw data throughput is high. Existing CSI-based solutions typically require high-performance network interface cards (NICs) and external processing power. There is a lack of optimized, localized frameworks that allow low-cost IoT microcontrollers like the ESP32 to capture and filter PHY-layer data in real-time without exceeding their limited RAM and processing capabilities.
\end{enumerate}

\subsection{Goals and Objectives}
\subsubsection{Goal}
The primary goal of this research is to develop and evaluate a hybrid RSSI-CSI localization framework capable of identifying and locating stationary, non-cooperative hidden IoT devices within cluttered indoor environments using low-cost edge hardware.

\subsubsection{Specific Objectives}
To achieve this goal, the following objectives must be met:
\begin{enumerate}
    \item \textbf{Characterization of Extreme Multipath:} To analyze and quantify the impact of signal shadowing and destructive interference on RSSI and CSI stability within a semi-structured classroom environment at the University of San Carlos.
    \item \textbf{Design of an Adaptive Hybrid Filter:} To develop a mathematical fusion algorithm that dynamically weighs Received Signal Strength Indicator (RSSI) for coarse range estimation and Channel State Information (CSI) sub-carrier amplitudes for fine-grained spatial fingerprinting.
    \item \textbf{Implementation of a Passive Sniffer Mesh:} To architect a distributed network of stationary ESP32 edge nodes that can capture and transmit Physical Layer (PHY) data to a centralized laptop gateway without requiring active device handshaking.
    \item \textbf{Comparative Performance Evaluation:} To validate the proposed hybrid approach by benchmarking its Mean Absolute Error (MAE) and signal stability against traditional RSSI-only trilateration methods in a cluttered indoor testbed.
\end{enumerate}

\subsection{Significance of the Study}
The success of this research offers significant practical and technical contributions across several domains:

\begin{itemize}
    \item \textbf{Cross-Sector Security and Privacy:} The ability to localize non-cooperative devices has immediate applications in \textbf{Military and Defense} for detecting rogue IoT devices, in \textbf{Corporate Business} for preventing industrial espionage, and in \textbf{Personal Privacy} for sweeping sensitive areas like hotel rooms or private residences for illicit surveillance.

    \item \textbf{Health and Safety Monitoring:} Beyond simple localization, the mastery of CSI amplitude analysis serves as a foundation for \textbf{Medical and Smart-Home} applications, where fine-grained signal variations can be utilized to monitor sleeping patterns, respiratory rates, or sudden falls without the use of intrusive cameras.

    \item \textbf{Educational and Institutional Infrastructure:} For institutions like the University of San Carlos, this study provides a blueprint for low-cost asset tracking within concrete-heavy structures where GPS is unavailable, ensuring the security of laboratory equipment in complex indoor layouts.

    \item \textbf{Democratization of Signal Research:} This research demonstrates that high-level Physical Layer (PHY) analysis is achievable on a limited budget. By utilizing lower-end microcontrollers (ESP32) instead of specialized thousand-dollar network interface cards, this study lowers the barrier of entry for students and local startups to explore advanced signal processing and IoT security.
\end{itemize}


\subsection{Scope and Limitations}

\subsubsection{Scope}
This research focuses on the development, implementation, and evaluation of a localization framework for detecting non-cooperative, stationary hidden IoT devices.

\begin{itemize}
    \item \textbf{Hardware and Platform:} The system is centered on the \textbf{ESP32 microcontroller} as a passive sniffing edge node and a \textbf{laptop gateway} for centralized data aggregation, signal filtering, and coordinate calculation.\smallskip
    \item \textbf{Localization Methodology:} The framework utilizes a \textbf{Hybrid Adaptive Fusion Algorithm} that integrates coarse RSSI-based trilateration with fine-grained CSI amplitude sub-carrier filtering to mitigate the destructive effects of indoor multipath propagation.\smallskip
    \item \textbf{Non-Associative Signal Acquisition:} The research leverages \textbf{802.11 control and management frames} (e.g., ACK, RTS/CTS, and Probe Responses). By capturing the PHY-layer preamble of these mandatory protocol responses, the system extracts CSI and RSSI data without requiring the target device to be formally associated or authenticated with the sniffer mesh.\smallskip
    \item \textbf{Environmental Context:} Testing is restricted to the \textbf{USC-TC Bunzel Building} classrooms. This environment is selected specifically for its high concentration of reflective surfaces (concrete and metal), which represent the most challenging conditions for indoor Wi-Fi localization.\smallskip
    \item \textbf{Performance Metrics:} The framework is evaluated based on its \textbf{Mean Absolute Error (MAE)} in meters and its signal-processing resilience to stationary physical obstructions such as furniture, equipment, and interior walls.\smallskip
\end{itemize}

\subsubsection{Limitations}
\begin{itemize}
    \item \textbf{Static Target Constraints:} The study is strictly limited to \textbf{stationary target nodes}. It does not address the complex signal variations, such as Doppler shifts or fast-fading effects, associated with high-speed mobile robotics or vehicular networks.\smallskip
    \item \textbf{Signal Dimensionality:} To maintain real-time computational efficiency on low-cost hardware, the research focuses exclusively on \textbf{CSI amplitude sub-carriers}. Complex phase-shift analysis and Angle-of-Arrival (AoA) estimations are excluded due to the lack of hardware clock synchronization on standard ESP32 modules.\smallskip
    \item \textbf{Frequency and Protocol Limitation:} The system operates solely on the \textbf{2.4 GHz ISM band} using the 802.11n protocol. While 5 GHz and 6 GHz bands offer different propagation characteristics, they are excluded to ensure hardware homogeneity across the sniffer mesh.\smallskip
    \item \textbf{Uncontrolled RF Interference:} The study cannot eliminate or fully account for external interference from existing university-wide Wi-Fi infrastructure (e.g., Eduroam) or transient mobile hotspots, which introduce non-deterministic noise into the experimental data.\smallskip
    \item \textbf{Secondary IoT Features:} While critical for industrial deployment, \textbf{battery life optimization} and \textbf{end-to-end network encryption} are beyond the scope of this research, which prioritizes signal-layer localization accuracy and filtering robustness.\smallskip
    \item \textbf{Dynamic Obstructions:} While the system accounts for static environmental clutter, the impact of \textbf{continuous human movement} (crowded rooms) during the localization process is excluded from the primary accuracy benchmarks to maintain a stable experimental baseline.\smallskip
    \item \textbf{Identity and Intent Agnostic:} The framework is designed to localize the physical coordinates of an RF source based on signal propagation. It does not perform \textbf{packet inspection, SSID validation, or MAC-layer security analysis}; therefore, the detection of specific cyber-attack vectors, such as ``Evil Twin'' or ``Man-in-the-Middle'' scenarios, is outside the scope of this signal-analysis-focused study.\smallskip
\end{itemize}

\section{Review of Related Literature}

% \subsection{Introduction to Chapter 2}
This chapter examines the literature that most directly informs the present study on hybrid RSSI and CSI-based localization for hidden IoT devices in cluttered indoor environments. Instead of presenting prior work as a disconnected inventory of methods, the review evaluates how recent scholarship addresses four interrelated concerns: non-cooperative device localization, the instability of RSSI under indoor multipath, the precision advantage of CSI-based fingerprinting, and the practical constraints of low-cost edge deployment. This analytical structure is important because the present study does not merely seek another indoor positioning system; it addresses a narrower and more difficult problem in which the target device is hidden, stationary, and potentially non-associative, while the sensing hardware must remain affordable. Accordingly, the discussion emphasizes where prior work is strong, where it remains limited, and how those limitations justify the proposed hybrid adaptive filtering framework.

\subsection{Related Literature}

\subsubsection{Hidden IoT Detection and Non-Cooperative Localization}
Recent scholarship has clarified that hidden-device localization is not simply a variant of generic indoor positioning; rather, it is a security-oriented problem defined by missing cooperation, incomplete prior knowledge, and the need for passive observation. Sharma et al. (2022), through Lumos, significantly advanced this discussion by demonstrating that diverse hidden IoT devices can be identified and localized in unfamiliar environments without assuming a conventional user-tracking setup. This contribution is important because it moves the field away from cooperative localization assumptions and toward adversarially realistic conditions. However, Lumos also reveals a recurring limitation in the literature: systems designed for hidden-device discovery often prioritize detection generality over fine-grained localization using low-cost hardware. In other words, they establish the security relevance of the problem but do not fully solve the engineering problem addressed by the present study.

Ju et al. (2025) strengthen this line of work by showing that channel fingerprints derived from RTS/CTS-triggered RSSI sequences can support the localization of hidden IoT devices in unknown environments. Their study is especially relevant because it validates RF-side signatures when direct device cooperation is unavailable. Yet the broader implication of this work is more important than the result itself: it suggests that the hidden-device problem cannot be solved reliably through network-layer assumptions alone. Even so, both Sharma et al. (2022) and Ju et al. (2025) leave unresolved the question of how such localization can be implemented in multipath-heavy educational interiors using commodity edge hardware. Thus, the literature establishes the problem setting convincingly, but it remains less conclusive on how hidden-device localization can be made simultaneously passive, low-cost, and precise under dense indoor reflections.

\subsubsection{RSSI-Based Indoor Localization Under Multipath Conditions}
RSSI remains attractive in the literature because it is universally available, lightweight to process, and easy to integrate into embedded systems. Nevertheless, recent studies converge on the same critical conclusion: these advantages are offset by weak spatial reliability in cluttered indoor environments. Abdel Ghany et al.\ (2020) compared narrowband LoRa CSI and RSSI in an outdoor IoT measurement campaign and reported stronger robustness for CSI-derived metrics than for RSSI under spatial and temporal variability \cite{ghany2020robustness}. For indoor Wi-Fi, Yang et al.\ (2021) surveyed recent localization scenarios and noted that RSSI-based methods often persist because of deployment convenience rather than peak precision \cite{yang2021survey}, complementing the earlier RSSI--CSI taxonomy of Yang et al.\ (2013) \cite{yang2013rssi2csi}. Leitch et al.\ (2023) reached a comparable conclusion by situating Wi-Fi RSSI among other indoor localization technologies and showing that its practical accessibility does not translate into resilience against indoor interference \cite{leitch2023indoor}.

What makes these studies especially relevant to the present work is not merely their critique of RSSI accuracy, but the implicit design lesson that follows from it. If RSSI is judged only by desk-level precision, it appears inadequate; however, if it is evaluated as a coarse and computationally efficient indicator of spatial proximity, its role becomes strategically useful. This distinction is often underdeveloped in prior studies, which tend to compare RSSI and CSI as competing features rather than as complementary ones. The present study departs from that binary framing. Instead of discarding RSSI because of its multipath sensitivity, it treats RSSI as an efficient first-stage estimator that constrains the search area before finer CSI analysis is invoked. The gap, therefore, is not simply that RSSI is inaccurate, but that prior literature insufficiently theorizes its value within adaptive hybrid pipelines.

\subsubsection{CSI-Based Fine-Grained Localization and Channel Fingerprinting}
In contrast to RSSI, CSI is increasingly treated as the more analytically powerful feature for indoor localization because it preserves sub-carrier-level channel structure and is therefore more responsive to subtle spatial differences. Li et al. (2024) showed that CSI feature fusion can improve localization performance while also supporting anomaly-aware processing, indicating that CSI is valuable not only for precision but also for robustness. Suroso et al.\ (2023) likewise emphasized that CSI-based methods are particularly suited to multipath-heavy environments because they exploit, rather than merely suffer from, rich channel variation \cite{suroso2023csi}. Reyes et al.\ (2025) extended this perspective by arguing that temporal stability is central to CSI fingerprinting performance; rich channel descriptors are only useful if they remain sufficiently consistent for reliable matching over time \cite{reyes2025temporal}.

However, the literature is more persuasive in demonstrating CSI's potential than in resolving its deployment cost. Many studies implicitly assume hardware stability, processing headroom, or calibration routines that are difficult to guarantee on low-cost microcontrollers. This is precisely where the present study positions its contribution. The issue is no longer whether CSI can outperform RSSI in theory; recent literature already suggests that it can. The unresolved issue is whether CSI can be operationalized in a passive, edge-feasible localization framework for hidden devices without inheriting the computational burden of research-grade platforms. Thus, CSI literature supports the study's fine-localization strategy, but it also sharpens the engineering gap that the proposed hybrid adaptive filter must address.

\subsubsection{Hybrid and Low-Cost Edge Localization Architectures}
The tension identified in the preceding literature naturally leads to hybrid and low-cost architectural research. If RSSI is efficient but imprecise, and CSI is precise but computationally heavier, then hybridization becomes less a matter of convenience and more a structural necessity. Mottakin et al.\ (2024), through SWiLoc, showed that ESP32-based Wi-Fi localization can be meaningfully implemented on commodity hardware, thereby challenging the assumption that effective indoor localization must depend on expensive research hardware \cite{mottakin2024swiloc}. Yet this work also underscores a critical limitation: feasibility at the hardware level does not automatically imply suitability for hidden, non-cooperative, multipath-heavy security scenarios. Put differently, low-cost localization has been demonstrated, but low-cost hidden-device localization under severe indoor clutter remains underexplored.

This is where hybrid architectures become analytically important. Li et al. (2024) imply that fusion is not merely an accuracy booster but a mechanism for making richer channel information usable under practical constraints. Even so, the current literature rarely specifies how role separation between RSSI and CSI should be designed in edge-constrained systems. Many hybrid approaches remain descriptive at the feature level but are less explicit about computational scheduling, fallback behavior, and passive operation. The present study addresses this omission by conceptualizing hybridization as an adaptive division of labor: RSSI provides coarse search-area reduction, while CSI is reserved for fine-point localization when higher spatial discrimination is required. This move from feature combination to role-based adaptive processing marks a more explicit response to the implementation gap left by prior work.

\subsection{Related Studies}
Recent empirical studies reveal a clear methodological shift: the field is moving away from proving that indoor localization is possible and toward determining under what constraints it remains reliable. Abdel Ghany et al.\ (2020) are important in this regard because their outdoor LoRa comparison of narrowband CSI and RSSI does more than rank two features; it empirically exposes how environmental noise reshapes the usefulness of each signal representation \cite{ghany2020robustness}. Li et al. (2024) push the literature further by showing that CSI feature fusion can improve localization even under anomalous or imperfect conditions, suggesting that robustness now depends on how features are combined rather than on any single feature alone. Ju et al. (2025) contribute by demonstrating that hidden-device localization is feasible in unknown environments, but their results also imply that the challenge is not simply identification accuracy; it is maintaining localization reliability when the target is passive or concealed.

The same empirical trend is visible in low-cost implementations. Mottakin et al.\ (2024) demonstrate that ESP32-based localization is achievable, while Reyes et al.\ (2025) show that temporal stability remains a major determinant of CSI fingerprinting accuracy \cite{mottakin2024swiloc,reyes2025temporal}. Considered together, these studies suggest that recent progress has been substantial but uneven: one line of work improves signal richness, another improves affordability, and another improves applicability to hidden devices. What remains insufficiently addressed is the integrated case in which all three conditions must be satisfied simultaneously. This is the empirical gap that gives the present study its relevance.

\subsection{Related Systems}
At the systems level, the literature makes clear that design assumptions strongly determine what kind of localization problem can actually be solved. Lumos (Sharma et al., 2022) is influential because it treats hidden-device localization as an operational system challenge rather than a purely algorithmic exercise. However, its broader scope also highlights a limitation for the present study: a system optimized for diverse hidden devices in unfamiliar environments is not necessarily optimized for fine-grained desk-level localization using inexpensive Wi-Fi edge nodes in a fixed educational setting. Thus, Lumos is best understood as a strong systems precedent, but not a direct architectural solution to the present study.

MotionCompass (He et al., 2021) sharpens this contrast further. By relying on movement-direction analysis for hidden wireless camera localization, it demonstrates that specialized systems can outperform more generic positioning pipelines when target behavior is exploitable. Yet this strength is also its limitation from the standpoint of the present research. The current study assumes stationary hidden devices, meaning that motion-derived cues cannot be treated as a reliable source of discrimination. Consequently, MotionCompass helps define the boundary of applicability for behavior-sensitive systems and, by contrast, underscores the need for a signal-centric framework grounded in passive RSSI and CSI measurements.

SWiLoc and related ESP32-based systems are valuable because they establish that low-cost embedded localization is feasible. Nevertheless, feasibility alone is not the same as adequacy. These systems are generally not designed around hidden, non-cooperative targets in severe indoor multipath, nor do they explicitly address adaptive coordination between coarse and fine localization stages. The present study therefore extends the systems literature not by rejecting these precedents, but by combining their affordability lessons with a security-driven localization objective and a more explicit hybrid processing logic.

\subsection{Synthesis of the Review}
The review makes four critical points clear. First, the literature consistently shows that RSSI remains operationally attractive but analytically insufficient for fine-grained indoor localization under severe multipath. Second, CSI is repeatedly validated as the more discriminative signal representation, yet its practical use is constrained by computational overhead, calibration demands, and hardware assumptions that do not always translate well to commodity edge platforms. Third, hidden-device localization research has matured enough to establish the security relevance of the problem, but many system designs still assume conditions that differ from the present study, such as target motion, richer sensing stacks, or broader detection goals that do not prioritize desk-level positioning. Fourth, low-cost localization research has improved accessibility, but it has not fully integrated hidden-device targeting, passive sensing, hybrid feature role separation, and cluttered indoor multipath into a single coherent framework.

These unresolved tensions define the research gap more precisely than a simple claim of ``limited prior work.'' The gap is not that RSSI, CSI, hidden-device systems, or ESP32 localization have been neglected individually. Rather, the gap lies in the absence of a framework that combines them in a way that is both methodologically coherent and practically deployable. Existing studies tend to trade one requirement for another: RSSI-based methods gain efficiency but lose precision; CSI-based methods gain precision but increase computational demand; hidden-device systems improve realism but often depend on assumptions not suitable for static classroom targets; and low-cost systems improve affordability without fully addressing passive, fine-grained localization under heavy multipath. The present study is justified precisely because it responds to this fragmented state of the literature through a hybrid adaptive filter in which RSSI performs coarse localization, CSI performs fine-point discrimination, and ESP32-laptop coordination makes the approach viable within a realistic academic environment. In this way, the study does not merely extend existing literature; it attempts to reconcile several strands of research that remain only partially connected in current scholarship.



\section{Methodology}

This chapter presents the methodology of the study. Based on Chapters 1 and 2, the present research is designed as a \textit{comparative analysis} between a \textbf{baseline} RSSI+CSI localization system reconstructed from related literature and a \textbf{proposed} hybrid localization system developed for this study. The comparison is necessary because the problem established in Chapter 1 is not simply whether hidden IoT devices can be localized, but whether a low-cost passive system can localize \textit{stationary, non-cooperative} devices in cluttered indoor environments more effectively than a conventional signal-fusion approach.

From Chapter 1, three methodological needs are clear. First, RSSI alone is unreliable in severe indoor multipath. Second, many accurate localization systems depend on motion, active cooperation, or expensive sensing hardware, making them unsuitable for hidden stationary devices. Third, CSI-based localization is promising but computationally heavier and less straightforward to deploy on commodity edge platforms. Chapter 2 refines this into a specific research gap: prior work validates RSSI, CSI, hidden-device localization, and low-cost ESP32 systems separately, but the literature does not fully integrate passive hidden-device targeting, low-cost deployment, and explicit role separation between coarse and fine localization in a single system \cite{yang2013rssi2csi,li2024efficient,mottakin2024swiloc,reyes2025temporal,ren2022lumos,ju2025hidden}. Accordingly, this chapter develops a methodology that compares a literature-grounded baseline against a proposed system intended to address that gap.

\subsection{Conceptual Framework}

The study uses an Input--Process--Output (IPO) framework to represent both systems. This framework is appropriate because both the baseline and proposed systems transform passive wireless observations into measurable localization outputs. The key difference lies in the structure of the process stage: the baseline performs direct localization from fused RSSI and CSI features, whereas the proposed system uses RSSI for coarse-zone estimation and CSI for fine refinement within the reduced search region.

\subsubsection{Baseline System IPO}

The baseline system is based on the signal-fusion direction described in the related literature, particularly the RSSI--CSI localization perspective discussed by Yang et al.\ (2013) and the CSI feature-fusion approach reported by Li et al.\ (2024) \cite{yang2013rssi2csi,li2024efficient}. Since the exact algorithmic steps are not fully enumerated in Chapter 2, the baseline is reconstructed as a realistic and literature-consistent RSSI+CSI localization pipeline in which both signal types are collected, preprocessed, fused, and used directly for position estimation. The conceptual structure of the baseline system is shown in Fig.~\ref{fig:baseline-ipo}.

\begin{figure}[H]
\centering
\includegraphics[width=\columnwidth]{prism-uploads/fig1-thesis.png}
\caption{Baseline IPO for the direct RSSI+CSI localization pipeline from related literature.}
\label{fig:baseline-ipo}
\end{figure}

\subsubsection{Proposed System IPO}

The proposed system uses the same general input sources but modifies the processing logic. Rather than fusing RSSI and CSI in a single direct localization stage, the proposed system assigns them distinct roles. RSSI is used for low-cost coarse spatial filtering, while CSI is used only after the candidate search space has been reduced. This directly addresses the research gap identified in Chapter 2, where hybrid approaches are discussed but not always operationalized as a staged decision framework under low-cost deployment constraints \cite{li2024efficient,mottakin2024swiloc,reyes2025temporal}. The conceptual structure of the proposed system is shown in Fig.~\ref{fig:proposed-ipo}.

\begin{figure}[H]
\centering
\includegraphics[width=\columnwidth]{prism-uploads/fig2-thesis.png}
\caption{Proposed IPO for the coarse-to-fine hybrid RSSI--CSI localization pipeline.}
\label{fig:proposed-ipo}
\end{figure}

\subsection{Baseline System (From Related Studies)}

\subsubsection{Algorithm Flow}

The baseline system is reconstructed from Chapter 2 studies that use RSSI and CSI as localization evidence. In particular, Yang et al.\ (2013) discuss RSSI and CSI as alternative and comparable wireless features for indoor positioning, while Li et al.\ (2024) show that CSI feature fusion improves localization robustness under imperfect conditions \cite{yang2013rssi2csi,li2024efficient}. Based on these studies, the baseline for the present work is defined as a \textit{direct RSSI+CSI localization pipeline} in which both RSSI and CSI are collected and fused before one final localization stage. This baseline is appropriate because it uses both RSSI and CSI as required, remains realistic, and reflects a standard fused-feature localization design.

In the baseline system, multiple ESP32 anchors are deployed around the test area and passively observe packets emitted by the hidden target device. Each anchor records packet timestamp, anchor identifier, RSSI value, and CSI amplitude information. These observations are sent to the laptop gateway, where invalid or incomplete records are discarded. The remaining records are grouped into short capture windows and aligned by timestamp so that observations from multiple anchors correspond to the same transmission interval.

After synchronization, the RSSI values are smoothed to reduce packet-to-packet fluctuations caused by fading and short-term interference. The CSI amplitudes are normalized to reduce scale variation across anchors and to preserve comparable subcarrier patterns. A fused feature vector is then formed from the processed RSSI and CSI observations. In the baseline system, this fused feature vector is used directly for localization, without separating coarse and fine inference stages. The gateway then returns an estimated position for the hidden device and records timing and reliability information for later comparison.

A realistic step-by-step logic for the baseline system is as follows:

\begin{enumerate}
    \item Initialize the ESP32 anchors and the laptop gateway.
    \item Passively capture packets from the hidden target device.
    \item Extract timestamp, anchor identifier, RSSI, and CSI amplitude features from each received packet.
    \item Send all observations to the gateway.
    \item Remove corrupt, incomplete, or duplicate records.
    \item Group records into synchronized capture windows.
    \item Smooth RSSI values over the window.
    \item Normalize CSI amplitudes and derive CSI feature vectors.
    \item Fuse RSSI and CSI features into a single localization input.
    \item Apply direct RSSI+CSI localization estimation.
    \item Output the estimated position and log processing metrics.
\end{enumerate}

The baseline pseudocode is shown below.

\begin{algorithm}[!ht]
\caption{Baseline Direct RSSI+CSI Localization}
\begin{algorithmic}[1]
\State Initialize anchors and gateway
\While{system is active}
    \State Capture packets from passive anchors
    \State Extract tuple $(id, t, RSSI, CSI, anchor)$
    \State Send tuples to gateway
    \State Remove corrupt or incomplete tuples
    \State Group tuples by capture window and target identifier
    \If{valid anchor observations $<$ minimum required}
        \State Mark window as insufficient and continue
    \EndIf
    \State Synchronize tuples across anchors
    \State Smooth RSSI values
    \State Normalize CSI amplitudes and derive feature vectors
    \State Fuse RSSI and CSI features
    \State Estimate position using direct RSSI+CSI localization
    \State Log position, latency, and packet sufficiency indicators
\EndWhile
\end{algorithmic}
\end{algorithm}
\vfill

\subsubsection{Algorithm Visualization}

Figure~\ref{fig:baseline-flow} summarizes the baseline direct RSSI+CSI localization workflow reconstructed from the related literature.

\begin{figure}[H]
\centering
\includegraphics[width=\columnwidth,height=0.70\textheight,keepaspectratio]{prism-uploads/fig3-thesis.png}
\caption{Baseline flowchart for the direct RSSI+CSI localization pipeline from related literature.}
\label{fig:baseline-flow}
\end{figure}

\subsection{Proposed System (This Study)}

\subsubsection{Algorithm Flow}

The proposed system improves on the baseline by changing how RSSI and CSI are used. Instead of treating both as equal inputs to a single direct localization step, the proposed method separates their roles. RSSI is used first to estimate a coarse region, while CSI is used second to refine the estimate only within that reduced candidate region. This design is grounded in the findings of Chapter 2: RSSI remains computationally simple but spatially unstable, whereas CSI is more discriminative but more computationally demanding \cite{yang2013rssi2csi,ghany2020robustness,li2024efficient,reyes2025temporal}. The proposed method therefore treats hybridization as \textit{adaptive sequencing} rather than mere feature fusion.

In the proposed system, the same set of ESP32 anchors passively collects RSSI and CSI observations from the stationary hidden device. The laptop gateway validates and synchronizes these records so that both the baseline and proposed systems begin from equivalent capture conditions. RSSI values are then used to estimate a coarse search area such as a zone, desk cluster, or bounded region of the room. This first stage reduces the search space and serves as a low-cost filter. The system then checks whether CSI observations are sufficient and stable for refinement. If they are, CSI features are compared only within the candidate positions belonging to the coarse zone. A refined estimate is then generated. If CSI observations are insufficient, the system falls back to the RSSI-based estimate instead of dropping the capture window completely.

This proposed workflow differs from the baseline in four key respects. First, it explicitly separates coarse and fine inference. Second, it reduces CSI computation by restricting refinement to a smaller candidate set. Third, it introduces fallback behavior for incomplete CSI conditions. Fourth, it better aligns with the objectives of the study, particularly low-cost deployment, robustness under imperfect measurements, and improved hidden-device localization in cluttered classrooms \cite{mottakin2024swiloc,li2024efficient,reyes2025temporal}.

The proposed algorithm follows these steps:

\begin{enumerate}
    \item Initialize the ESP32 anchors and the laptop gateway.
    \item Passively capture packets from the hidden target device.
    \item Extract timestamp, anchor identifier, RSSI, and CSI from each valid packet.
    \item Remove corrupted or incomplete records.
    \item Group observations into synchronized capture windows.
    \item Smooth RSSI values and normalize CSI amplitudes.
    \item Use RSSI values to estimate a coarse search zone.
    \item Check whether CSI observations are sufficient and stable.
    \item If CSI is sufficient, evaluate candidate positions only within the coarse zone using CSI features.
    \item Generate the refined hybrid position estimate.
    \item If CSI is insufficient, output the RSSI-based fallback estimate.
    \item Log final estimate, mode used, latency, and reliability indicators.
\end{enumerate}

\newpage
The proposed pseudocode is shown below.

\begin{algorithm}[!ht]
\caption{Proposed Hybrid Coarse-to-Fine RSSI--CSI Localization}
\begin{algorithmic}[1]
\State Initialize anchors and gateway
\While{system is active}
    \State Capture packets from passive anchors
    \State Extract tuple $(id, t, RSSI, CSI, anchor)$
    \State Send tuples to gateway
    \State Remove corrupt or incomplete tuples
    \State Group tuples by capture window and target identifier
    \If{valid anchor observations $<$ minimum required}
        \State Mark window as insufficient and continue
    \EndIf
    \State Synchronize tuples across anchors
    \State Smooth RSSI values
    \State Normalize CSI amplitudes
    \State Estimate coarse zone using RSSI
    \If{CSI observations are sufficient and stable}
        \State Evaluate candidate positions within the coarse zone using CSI
        \State Generate refined hybrid estimate
    \Else
        \State Use RSSI coarse estimate as fallback
    \EndIf
    \State Log final position, mode used, latency, and packet sufficiency indicators
\EndWhile
\end{algorithmic}
\end{algorithm}

\vfill
\newpage

\subsubsection{Algorithm Visualization}



Figure~\ref{fig:proposed-flow} presents the proposed coarse-to-fine hybrid workflow with fallback behavior when CSI observations are insufficient.

\begin{figure}[H]
\centering
\includegraphics[width=\columnwidth,height=0.70\textheight,keepaspectratio]{prism-uploads/fig4-thesis.png}
\caption{Proposed flowchart for the coarse-to-fine hybrid RSSI--CSI localization pipeline with fallback.}
\label{fig:proposed-flow}
\end{figure}

\subsection{Comparison / Analysis}

The study compares the baseline and proposed systems under identical operating conditions. Both systems use the same hidden target device, the same anchor positions, the same classroom environment, the same capture duration, the same synchronized windows, and the same ground-truth coordinates. This is necessary to ensure that any performance difference is due to the localization method rather than to unequal measurement conditions.

The comparison is based on the following metrics:

\begin{enumerate}
    \item \textbf{Localization accuracy:} measured using Mean Absolute Error (MAE), Root Mean Square Error (RMSE), median error, and 90th percentile error relative to known target coordinates.
    \item \textbf{Response time / latency:} measured from the time a valid capture window is formed to the time a position estimate is produced.
    \item \textbf{Stability / noise resistance:} measured through repeated-trial estimate spread, variance at fixed positions, and performance under degraded sensing conditions.
    \item \textbf{Efficiency:} measured through usable window rate, localization availability, and computational processing cost per window.
\end{enumerate}

The baseline is expected to provide a simpler direct localization pipeline, but it may be more sensitive to noisy fused observations because RSSI and CSI are combined without explicit search-space reduction. The proposed system is expected to improve robustness and fine-grained accuracy by using RSSI for search-area reduction before invoking CSI refinement. However, this expectation will be validated empirically rather than assumed.

\subsubsection{Experimental Environment (TechHub Floor Plan)}

The experimental environment is the selected TechHub classroom or laboratory test area where the hidden device is placed at predetermined static positions. Multiple ESP32 anchors are positioned around the perimeter or corners of the room to provide spatial diversity. The laptop gateway is placed at a monitoring station where all observations are synchronized and processed. Since the study concerns stationary hidden-device localization rather than continuous user tracking, the main experimental concern is repeated placement of the target at known positions rather than long movement trajectories. If movement is included for packet diversity, it will be treated as controlled repositioning between trials rather than continuous tracking.

Figure~\ref{fig:techhub-floorplan} shows the uploaded TechHub floor plan used as the experimental layout reference.

\begin{figure}[H]
\centering
\includegraphics[width=\columnwidth]{prism-uploads/techhub.jpeg}
\caption{Uploaded TechHub floor plan used as the experimental layout reference for anchor placement, gateway location, and target test positions.}
\label{fig:techhub-floorplan}
\end{figure}

In this layout, \texttt{A1--A4} denote the ESP32 anchors, \texttt{P1--P10} denote predetermined hidden-device test positions, and the gateway is the laptop monitoring station. The signal sources in the experiment are the target device transmissions observed by the anchors. Because the target is stationary during each run, each point is tested through repeated capture windows rather than a continuous movement path.

\subsection{Systems Analysis and Design}

\subsubsection{Functional Requirements}

The system must satisfy the following functional requirements:

\begin{enumerate}
    \item It must passively capture RSSI and CSI-related packet observations from a stationary hidden IoT device.
    \item It must support multi-anchor data collection using ESP32 nodes.
    \item It must synchronize observations across anchors at the gateway.
    \item It must support execution of both the baseline and proposed localization systems.
    \item It must generate a location estimate for each valid capture window.
    \item It must log all outputs needed for comparative analysis.
    \item It must support repeated trials across multiple known target positions.
\end{enumerate}

\subsubsection{Non-Functional Requirements}

The system must also satisfy the following non-functional requirements:

\begin{enumerate}
    \item \textbf{Low cost:} the platform must remain deployable using commodity ESP32 anchors and a laptop gateway.
    \item \textbf{Reproducibility:} the capture and processing workflow must remain consistent across repeated trials.
    \item \textbf{Robustness:} the system must tolerate packet noise, multipath distortion, and partial data loss.
    \item \textbf{Efficiency:} the gateway must process synchronized windows within acceptable delay.
    \item \textbf{Maintainability:} capture, preprocessing, estimation, and logging modules must be separable and modifiable.
\end{enumerate}

\subsubsection{System Architecture}

The system architecture follows a distributed-sensing and centralized-processing design. Multiple ESP32 anchors perform passive packet observation, while the laptop gateway performs synchronization, preprocessing, localization, and logging. The architecture contains four layers:

\begin{enumerate}
    \item \textbf{Target layer:} the hidden IoT device emitting observable wireless packets.
    \item \textbf{Anchor layer:} ESP32 nodes capturing RSSI and CSI-related measurements.
    \item \textbf{Gateway layer:} laptop executing either the baseline or proposed localization pipeline.
    \item \textbf{Evaluation layer:} the analysis stage where outputs are compared against ground truth.
\end{enumerate}

This architecture is consistent with Chapter 2, particularly the literature showing that low-cost ESP32 localization is feasible and that more computationally expensive signal interpretation is better handled centrally \cite{mottakin2024swiloc,li2024efficient}.

\subsection{System Components}

\subsubsection{Hardware}

The hardware components of the system are:

\begin{enumerate}
    \item ESP32 microcontrollers serving as passive anchor nodes;
    \item laptop serving as the gateway and processing unit;
    \item the hidden IoT target device used during localization trials;
    \item power supplies and mounting fixtures for stable anchor placement; and
    \item the TechHub classroom test environment containing walls, desks, and other multipath-causing structures.
\end{enumerate}

Each ESP32 anchor is responsible for local observation of target packets. The laptop is responsible for centralized synchronization, inference, and logging. The classroom environment is part of the method because the research specifically examines cluttered indoor multipath conditions.

\subsubsection{Software}

The software components are:

\begin{enumerate}
    \item anchor firmware for passive packet capture and feature extraction;
    \item gateway synchronization and preprocessing software;
    \item baseline direct RSSI+CSI localization module;
    \item proposed coarse-to-fine hybrid localization module;
    \item logging and metric computation module; and
    \item analysis routines for comparing baseline and proposed outputs.
\end{enumerate}

\subsubsection{Network}

The network arrangement is a local passive observation network. The hidden target device transmits packets that are observed by multiple ESP32 anchors. The anchors forward packet features or aggregated window data to the laptop gateway through the local network. The target device does not need to authenticate, associate, or cooperate with the system. This is important because the study explicitly concerns hidden and non-cooperative devices, consistent with the systems motivation established in Chapter 2 \cite{ren2022lumos,ju2025hidden}.

\subsection{Testing and Validation Plan}

Testing will be conducted at three levels for both systems.

\textbf{Unit testing} will be applied to individual modules such as packet parsing, timestamp synchronization, RSSI smoothing, CSI normalization, feature fusion, coarse-zone estimation, CSI refinement, fallback logic, and metric computation. Each module will be tested independently using known inputs and expected outputs so that failures can be isolated before full deployment. Validation at this level will focus on input-output correctness, execution consistency, and error handling for incomplete or malformed records.

\textbf{Integration testing} will verify that anchor observations are correctly transferred to the gateway, that synchronized capture windows are formed properly, and that both baseline and proposed systems receive equivalent inputs from the same sensing process. This stage will specifically test the interaction between ESP32 anchors, the laptop gateway, the preprocessing pipeline, and the localization modules. Integration will be considered successful when packet records are delivered without structural mismatch, timestamps are aligned correctly, and the gateway generates valid localization windows from multi-anchor observations.

\textbf{System testing} will be conducted in the actual TechHub environment using predetermined hidden-device positions. For each position, multiple capture windows will be collected so that accuracy, stability, and latency can be measured across repeated trials. At this level, the full sensing, preprocessing, localization, and logging chain will be evaluated as one operational system under realistic environmental conditions.

Controlled testing conditions will include:

\begin{enumerate}
    \item repeated trials at each fixed target position;
    \item trials under normal classroom clutter;
    \item trials near reflective or obstructed areas;
    \item trials with reduced anchor availability; and
    \item trials in which CSI quality becomes insufficient, to test fallback behavior in the proposed system.
\end{enumerate}

Both systems will be evaluated using the same anchor arrangement, target positions, observation duration, and ground-truth references. This controlled design ensures fair and measurable comparison.

\subsubsection{Real-World (IRL) Testing}

Real-World (IRL) testing will be performed after unit, integration, and controlled system testing have confirmed that the hardware and software components are functioning correctly. The purpose of IRL testing is to determine whether the proposed localization system can maintain acceptable performance when deployed under realistic operating conditions in the TechHub environment, where signal blockage, human activity, device orientation changes, and irregular packet availability may affect the quality of RSSI and CSI observations. Unlike controlled tests that isolate variables as much as possible, IRL testing is intended to measure the practical behavior of the complete system under actual usage constraints.

The IRL testing procedure will follow the steps below.

\begin{enumerate}
    \item \textbf{Step 1: Site Preparation and Deployment Verification} \\
    \textit{Purpose:} To confirm that the physical deployment of the anchors and gateway matches the intended experimental layout before real-world trials begin. \\
    \textit{Procedure:} The ESP32 anchors will be mounted at their assigned positions in the TechHub environment, powered continuously, and connected to the laptop gateway. Anchor visibility, packet reception, timestamp consistency, and gateway logging will be checked before any localization trial is started. \\
    \textit{Expected Outcomes:} All anchors should report observations to the gateway without persistent communication failure. The gateway should correctly register anchor identity, packet records, and synchronized windows. \\
    \textit{Metrics and Validation:} Anchor availability rate, packet reception continuity, gateway connectivity status, and synchronization success rate will be recorded. Deployment will be accepted only if all required anchors remain operational for the full calibration interval.

    \item \textbf{Step 2: Baseline Environmental Capture} \\
    \textit{Purpose:} To characterize the normal RF condition of the TechHub site before target trials are executed. \\
    \textit{Procedure:} With the anchors active and the target device absent or inactive, the system will collect background packet observations over a fixed time interval. Environmental notes such as room occupancy, open or closed doors, and nearby active wireless devices will also be recorded. \\
    \textit{Expected Outcomes:} The system should produce a baseline record of ambient packet activity and noise conditions that can later be used to interpret changes observed during localization trials. \\
    \textit{Metrics and Validation:} Background packet rate, signal variability, observation gaps, and any anomalous interference events will be logged. These records will serve as a reference for determining whether later localization failures are caused by environmental noise rather than by algorithmic error.

    \item \textbf{Step 3: Real-Position Localization Trials} \\
    \textit{Purpose:} To evaluate localization accuracy under realistic deployment conditions using actual hidden-device placements. \\
    \textit{Procedure:} The hidden IoT target device will be placed at predetermined real positions in the TechHub environment, including open areas, partially obstructed areas, and positions near reflective surfaces. At each position, the system will collect multiple capture windows while both the baseline and proposed localization pipelines process the same observation conditions. \\
    \textit{Expected Outcomes:} Both systems should return estimable positions for most valid windows, while the proposed system is expected to achieve lower error and more stable estimates than the baseline. \\
    \textit{Metrics and Validation:} Mean Absolute Error (MAE), Root Mean Square Error (RMSE), median error, 90th percentile error, and coarse-zone accuracy will be computed against known ground-truth coordinates. Validation will require that comparisons be made using identical target positions and anchor layouts for both systems.

    \item \textbf{Step 4: Performance and Latency Evaluation Under Live Operation} \\
    \textit{Purpose:} To determine whether the system remains responsive enough for practical use during actual operation. \\
    \textit{Procedure:} During active localization trials, the timestamp at which a valid synchronized window is formed and the timestamp at which the final estimate is produced will be recorded automatically by the gateway. The same measurement process will be applied to both baseline and proposed systems. \\
    \textit{Expected Outcomes:} The system should produce estimates within an acceptable delay range for repeated hidden-device detection trials. The proposed system may require more computation than the baseline, but the increase should remain operationally manageable. \\
    \textit{Metrics and Validation:} End-to-end latency, processing time per window, usable estimate rate, and localization availability rate will be measured. Validation will compare whether any latency increase introduced by the proposed method is justified by corresponding gains in accuracy and robustness.

    \item \textbf{Step 5: Reliability and Robustness Testing in Real Use Conditions} \\
    \textit{Purpose:} To test whether the system remains functional when realistic disturbances affect sensing quality. \\
    \textit{Procedure:} Additional IRL trials will be conducted while introducing practical disturbances such as partial human movement in the environment, one temporarily unavailable anchor, minor device orientation changes, and windows with weak or incomplete CSI observations. These conditions will not be artificially exaggerated beyond what may normally occur in the TechHub setting. \\
    \textit{Expected Outcomes:} The baseline system may show greater instability under these conditions, while the proposed system should maintain better continuity through coarse-zone filtering and fallback behavior. \\
    \textit{Metrics and Validation:} Estimate stability, fallback success rate, anchor reduction tolerance, packet sufficiency rate, and successful estimate rate under disturbance will be recorded. Reliability will be validated by comparing how often each system continues to return usable estimates without severe error escalation.

    \item \textbf{Step 6: Hardware and Software Operational Validation} \\
    \textit{Purpose:} To verify that both the hardware platform and software pipeline remain dependable during extended IRL use. \\
    \textit{Procedure:} Throughout the real-world trials, the system will monitor anchor uptime, gateway logging continuity, packet parsing integrity, synchronization consistency, and correctness of recorded outputs. Any mismatch between observed device state and logged system state will be documented immediately. \\
    \textit{Expected Outcomes:} The ESP32 anchors should remain powered and responsive, the laptop gateway should complete logging without crash or data corruption, and the software pipeline should continue to process incoming windows without structural failure. \\
    \textit{Metrics and Validation:} Hardware uptime, dropped-record count, parsing error rate, synchronization failure rate, and log completeness will be measured. Validation at this step requires that the system finish each IRL session with complete and analyzable records.

    \item \textbf{Step 7: User-Operator Interaction and Result Review} \\
    \textit{Purpose:} To confirm that the deployed system can be operated consistently by the researcher during actual testing sessions. \\
    \textit{Procedure:} The operator will perform the standard workflow of positioning anchors, selecting test points, initiating capture, monitoring gateway output, and recording the resulting estimates. Any points at which the operator must intervene because of unclear output, inconsistent logs, or repeated failed captures will be documented. \\
    \textit{Expected Outcomes:} The system should remain understandable and manageable during use, with clear trial start, data capture, and result logging behavior. \\
    \textit{Metrics and Validation:} Number of interrupted trials, repeated trial initiations, operator-noted anomalies, and log interpretability will be documented. While the system is not designed as an end-user interface product, this step validates practical operability during experimentation.

    \item \textbf{Step 8: Final IRL Comparative Validation} \\
    \textit{Purpose:} To determine whether the proposed system demonstrates practical improvement over the baseline under actual TechHub conditions. \\
    \textit{Procedure:} After all IRL trials are completed, the recorded outputs from the baseline and proposed systems will be compared using the same set of target positions, trial counts, and environmental notes. Results from normal and disturbed conditions will be analyzed together to determine whether the proposed method offers consistent performance advantages. \\
    \textit{Expected Outcomes:} The proposed system should demonstrate measurable gains in localization quality, continuity, or robustness relative to the baseline direct RSSI+CSI method. \\
    \textit{Metrics and Validation:} Comparative MAE, RMSE, latency, stability, availability, and robustness indicators will be summarized. Validation will be achieved if the proposed system shows improved or more reliable localization performance without unacceptable operational cost under real-world conditions.
\end{enumerate}

\subsection{Validation and Deployment Plan}

Deployment begins with calibration of anchor positions inside the TechHub environment and verification that all ESP32 nodes can report observations to the laptop gateway. After communication and timestamp alignment are confirmed, baseline measurements of the indoor RF environment will be collected. Formal experimentation will then proceed by placing the hidden target device at predetermined coordinates and collecting repeated capture windows for both baseline and proposed processing.

User interaction during deployment is minimal. The operator mainly sets anchor positions, selects the target test point, initiates capture, and reads the resulting localization output from the gateway. Since the study is focused on hidden-device localization rather than an interactive end-user application, the main operational priority is consistency of deployment rather than complex user control.

Improvement over the baseline will be validated through direct comparison of localization error, latency, stability, and efficiency under identical conditions. The proposed system will be considered superior only if it demonstrates measurable gains in localization quality or robustness without introducing impractical processing overhead. In this way, the validation and deployment plan remains aligned with the main goal of the study: to determine whether a hybrid RSSI--CSI framework on low-cost hardware can outperform a literature-grounded RSSI+CSI baseline for stationary hidden-device localization in cluttered indoor environments.

\section*{IV. SUMMARY AND HYPOTHESIS}

\subsection*{Summary}

This study presents a low-cost indoor localization framework for identifying stationary, non-cooperative hidden IoT devices in cluttered indoor environments. The project addresses the limitations of conventional RSSI-based localization in spaces affected by severe multipath propagation, signal shadowing, and unstable indoor measurements. In response to this problem, the system is designed to use both Received Signal Strength Indicator (RSSI) and Channel State Information (CSI) in order to improve localization reliability and reduce error under realistic classroom conditions.

The proposed framework is implemented through a distributed architecture composed of ESP32 edge nodes that act as passive sniffers and a centralized laptop gateway that performs aggregation, preprocessing, synchronization, and localization estimation. At the system level, packet observations are passively captured from the hidden target device, validated, grouped into synchronized capture windows, and processed for localization. The literature-grounded baseline system applies direct fused localization by smoothing RSSI, normalizing CSI, and estimating position in a single-stage pipeline. In contrast, the proposed system adopts a coarse-to-fine hybrid strategy in which RSSI is first used for coarse zone estimation, after which CSI is applied only within the reduced candidate region for fine refinement. When CSI observations are insufficient or unstable, the proposed system retains functionality through an RSSI-based fallback estimate.

The study is structured as a comparative evaluation between the baseline and proposed systems under identical deployment conditions, anchor placements, target positions, and observation windows. Performance is assessed using measurable indicators such as Mean Absolute Error (MAE), Root Mean Square Error (RMSE), latency, estimate stability, availability, and robustness under disturbed conditions. Through this design, the project aims to determine whether a hybrid RSSI--CSI localization framework on commodity hardware can provide more accurate, stable, and operationally practical hidden-device localization than a direct fused baseline in cluttered indoor environments.

\subsection*{Hypothesis}

The study hypothesizes that the proposed coarse-to-fine hybrid RSSI--CSI localization framework will achieve significantly better localization performance than the literature-grounded baseline direct RSSI+CSI system for stationary hidden IoT devices in cluttered indoor environments, as measured by lower localization error, improved estimate stability, and better robustness under imperfect sensing conditions.

\noindent\textbf{Null Hypothesis (H0).} There is no significant difference between the proposed coarse-to-fine hybrid RSSI--CSI localization framework and the literature-grounded baseline direct RSSI+CSI system in terms of localization accuracy, stability, latency, or robustness for stationary hidden IoT devices in cluttered indoor environments.

\noindent\textbf{Alternative Hypothesis (H1).} The proposed coarse-to-fine hybrid RSSI--CSI localization framework provides significantly better localization performance than the literature-grounded baseline direct RSSI+CSI system, as evidenced by lower Mean Absolute Error (MAE) and Root Mean Square Error (RMSE), more stable estimates, and improved robustness under realistic indoor disturbances.

\nocite{ren2022lumos,ju2025hidden,ghany2020robustness,li2024efficient,yang2021survey,yang2013rssi2csi,wu2012csi,zaruba2007singleap,sensors2022indoor3d,xie2015urbancanyon,leitch2023indoor,suroso2023csi,mottakin2024swiloc,reyes2025temporal,he2021motioncompass}

\bibliographystyle{IEEEtran}
\bibliography{references}

\end{document}
